\documentclass[tikz,border=1cm]{standalone}
\usepackage{xcolor}
\usepackage{amsmath,amssymb}

\usepackage[utf8]{inputenc}
\usepackage[spanish]{babel}
\usepackage[
    letterpaper,
    top=1.5cm,
    bottom=2cm,
    left=2cm,
    right=2cm
]{geometry}

\usepackage{graphicx}
\usepackage{fancyhdr}
\usepackage[table,xcdraw]{xcolor}
\usepackage{tikz}
\usetikzlibrary{shapes.geometric, arrows, positioning, fit, backgrounds, calc, babel}
\usepackage{ifpdf}
\usepackage{blindtext}
\usepackage[cmex10]{amsmath}
\usepackage{amssymb}
\usepackage{amsfonts}
\usepackage{float}
\usepackage{siunitx}
\usepackage{listings}
\usepackage{hyperref}
\usepackage{subcaption}
\usepackage{booktabs}
\usepackage{url}
\usepackage{wrapfig}
\usepackage{titlesec}
\usepackage[export]{adjustbox}
\usepackage{imakeidx}
\usepackage{parskip}
\usepackage[numbers]{natbib}

\usepackage{ragged2e}
\usepackage{setspace}
\usepackage{array}
\usepackage{longtable}

% Definición de un nuevo tipo de columna para evitar "Underfull \hbox" en tablas
\newcolumntype{L}[1]{>{\RaggedRight\arraybackslash}p{#1}}

\hypersetup{colorlinks=false,bookmarksopen=true,linkbordercolor={1 1 1}}

\newcommand{\rfigura}[1]{Fig.\ref{#1}}
\newcommand{\rtabla}[1]{Tabla \ref{#1}}

% Datos del informe
\newcommand{\Curso}{IEE2913 --- Diseño Eléctrico (Capstone)}
\newcommand{\TituloInforme}{Informe de Avance 1 (5\%)}
\newcommand{\NumeroGrupo}{10}
\newcommand{\NombreProyecto}{SupaClock: Wearable Biométrico Modular}
\newcommand{\Integrantes}{Tomás Avendaño, Benjamín Sepúlveda, Pablo Uribe}
\newcommand{\FechaEntrega}{7 de abril de 2026}

\lstset{
  basicstyle=\ttfamily\small,
  breaklines=true,
  breakatwhitespace=true,
  postbreak=\mbox{\textcolor{red}{$\hookrightarrow$}\space},
  captionpos=b,
  frame=single,
  numbers=left,
  numberstyle=\tiny\color{gray},
  language=C,
  keywordstyle=\color{blue}\bfseries,
  commentstyle=\color{gray}\itshape,
  stringstyle=\color{red!70!black}
}


\begin{document}
\begin{tikzpicture}[
                node distance=1.5cm and 2cm,
                device/.style={rectangle, rounded corners, draw=black, thick, top color=white, bottom color=green!10, text width=3cm, align=center, minimum height=1.5cm, font=\small\bfseries},
                app/.style={rectangle, rounded corners, draw=black, thick, top color=white, bottom color=blue!10, text width=3cm, align=center, minimum height=1.5cm, font=\small\bfseries},
                cloud/.style={ellipse, draw=black, thick, top color=white, bottom color=orange!10, text width=3.5cm, align=center, minimum height=2cm, font=\small},
                db/.style={cylinder, shape border rotate=90, aspect=0.25, draw=black, thick, top color=white, bottom color=orange!20, text width=2.5cm, align=center, minimum height=2cm, font=\footnotesize},
                func/.style={rectangle, draw=black, dashed, fill=gray!10, text width=3.5cm, align=center, minimum height=1.5cm, font=\footnotesize},
                arrow/.style={thick, ->, >=stealth},
                double_arrow/.style={thick, <->, >=stealth},
                legend/.style={rectangle, draw=black!50, fill=white, inner sep=5pt, font=\scriptsize}
            ]%
            % Nodos Principales
            \node[device] (watch) {SupaClock\\(Nodo Edge)};
            \node[app, right=3.5cm of watch] (android) {Aplicación Móvil\\(Gateway Android)};
            %
            % Nodos de Nube
            \node[cloud, right=2.5cm of android] (firebase) {\textbf{Ecosistema Firebase}};
            \node[db, right=1cm of firebase, yshift=1.5cm] (firestore) {Cloud Firestore\\(Perfiles)};
            \node[func, right=1cm of firebase, yshift=-1.5cm] (functions) {\textbf{Cloud Functions}\\(Python)\\ \textit{Algoritmo ECG}};
            %
            % Nodo Raspberry (Futuro)
            \node[device, bottom color=purple!10, dashed, below=2cm of android] (rpi) {Raspberry Pi 3B+\\(Edge Server)};
            %
            % Conexiones BLE
            \draw[arrow, blue] ([yshift=0.5cm]watch.east) -- node[above, font=\scriptsize, text=black, align=center] {BLE (1 Hz)\\Telemetría Continua} ([yshift=0.5cm]android.west);
            \draw[arrow, red] ([yshift=-0.5cm]watch.east) -- node[below, font=\scriptsize, text=black, align=center] {BLE (100 Hz)\\ECG Crudo} ([yshift=-0.5cm]android.west);
            %
            % Conexiones App -> Nube
            \draw[double_arrow] (android) -- node[above, font=\scriptsize] {Wi-Fi / 5G} (firebase);
            %
            % Conexiones internas nube
            \draw[arrow] (firebase) -- (firestore);
            \draw[arrow] (firebase) -- (functions);
            \draw[arrow, dashed] (functions) -- node[right, font=\scriptsize] {Actualiza BPM} (firestore);
            %
            % Conexión Futura
            \draw[double_arrow, dashed, purple] (android) -- node[right, font=\scriptsize] {Proyección local} (rpi);
            %
            % === LEYENDA ===
            \node[legend, below=1cm of watch, xshift=3cm] (leg) {
                \begin{tabular}{cl}
                    \tikz\draw[fill=green!10] (0,0) rectangle (0.3,0.2);          & Dispositivo físico      \\
                    \tikz\draw[fill=blue!10] (0,0) rectangle (0.3,0.2);           & Aplicación de software  \\
                    \tikz\draw[fill=orange!10] (0,0) rectangle (0.3,0.2);         & Servicio en la nube     \\
                    \tikz\draw[fill=purple!10, dashed] (0,0) rectangle (0.3,0.2); & Proyección futura       \\
                    \tikz\draw[blue, thick, ->] (0,0.1) -- (0.4,0.1);             & BLE: Telemetría (1 Hz)  \\
                    \tikz\draw[red, thick, ->] (0,0.1) -- (0.4,0.1);              & BLE: ECG crudo (100 Hz) \\
                    \tikz\draw[purple, thick, dashed, <->] (0,0.1) -- (0.4,0.1);  & Enlace futuro           \\
                \end{tabular}
            };
        \end{tikzpicture}
\end{document}
